
\subsection{双模同态之间的同伦}
\label{subsec:5-3-3}
\liftIndex{I195}

本小节要证明下述结论。
\begin{theorem}{5.3.27}\label{thm:5-3-27}
定理 5.3.14 中的
$\mathfrak E$-弱滤过
$A_\infty$ 双模同态 $\phi$，
在双模同态的同伦意义下
不依赖于构造时所作的各种选择，
而只依赖于两族哈密顿同痕
$\{\psi^{(0)}_\rho\}_\rho$、
$\{\psi^{(1)}_\rho\}_\rho$
相对边界的同伦类。
\end{theorem}

\begin{proof}
为证明定理 5.3.27，
只须说明上一小节所构造的
$\mathfrak E$-弱滤过
$A_\infty$ 双模同态 $\phi$
在同伦意义下不依赖于各种选择。
定义 $\phi$ 时作了以下选择：
\begin{itemize}
\item 哈密顿同痕
$\psi^{(0)}=\{\psi^{(0)}_\rho\}_\rho$、
$\psi^{(1)}=\{\psi^{(1)}_\rho\}_\rho$；
\item 相容殆复结构的双参数族
$\{J_{\tau,t}\}_{\tau,t}$；
\item 多值截面 $\mathfrak s$；
\item $R'_{h'}$ 上光滑奇异单形的
可数集 $\mathcal X_3(R'_{h'})$。
\end{itemize}
为标明这些选择，
第一套选择所得的弱滤过同态记为
\[
\phi_{s=0}^{
\psi^{(0),s=0},\psi^{(1),s=0},
\{J_{\tau,t,s=0}\}_{\tau,t},
\mathfrak s_{s=0}}.
\]
它立于一对滤过 $A_\infty$ 代数同态
\[
\bigl(
\mathfrak f_{s=0}^{0,\psi^{(0),s=0},
\{J^{(0)}_\rho\}_\rho,
\mathfrak s_{s=0}^{\operatorname{top}(\rho)}},
\mathfrak f_{s=0}^{1,\psi^{(1),s=0},
\{J^{(1)}_\rho\}_\rho,
\mathfrak s_{s=0}^{\operatorname{top}(\rho)}}
\bigr)
\]
之上；
这里采用第 4.6.2 小节的记号。
第二套选择
$\psi^{(0),s=1}$、
$\psi^{(1),s=1}$、
$\{J_{\tau,t,s=1}\}_{\tau,t}$、
$\mathfrak s_{s=1}$
类似给出弱滤过双模同态
\[
\phi_{s=1}^{
\psi^{(0),s=1},\psi^{(1),s=1},
\{J_{\tau,t,s=1}\}_{\tau,t},
\mathfrak s_{s=1}},
\]
它立于相应的一对滤过
$A_\infty$ 代数同态之上。

假设
$\psi^{(0),s=1}$、
$\psi^{(1),s=1}$
分别与
$\psi^{(0),s=0}$、
$\psi^{(1),s=0}$
同伦。
具体地，取双参数族
$\{\psi^{(i)}_{\rho,s}\}_{\rho,s}$
（$i=0,1$），使
\[
\begin{cases}
\psi^{(i)}_{\rho,0}
 =\psi^{(i),s=0}_\rho,\qquad
\psi^{(i)}_{\rho,1}
 =\psi^{(i),s=1}_\rho,\qquad
\psi^{(i)}_{0,s}=\mathrm{id},\\
\psi^{(i)}_{1,s}(L^{(i)})
 =\psi^{(i),s=0}_1(L^{(i)})
 =\psi^{(i),s=1}_1(L^{(i)}).
\end{cases}
\]
置
\[
\psi^{(i)}_{1,s}(L^{(i)})
=L^{(i)\prime}
\qquad(i=0,1).
\]
对两个端点 $i=0,1$，
另定义
$\psi_{\rho,t,i}\colon M\to M$
（$0\le\rho\le1$）为
\[
\psi_{\rho,t,i}
=\psi^{(1)}_{\rho t,i}
\circ\psi^{(0)}_{\rho(1-t),i}.
\]

第 4.6.2 小节已构造
从
$\mathfrak f_{s=0}^{i,\psi^{(i),s=0},
\{J^{(i)}_\rho\}_\rho,
\mathfrak s_{s=0}^{\operatorname{top}(\rho)}}$
到
$\mathfrak f_{s=1}^{i,\psi^{(i),s=1},
\{J^{(i)}_\rho\}_\rho,
\mathfrak s_{s=1}^{\operatorname{top}(\rho)}}$
的同伦。
那里构造的同伦依赖以下数据
（$i=0,1$）：
\begin{itemize}
\item 哈密顿同痕
$\{\psi^{(i)}_{1,s}\}_s$，
满足
$\psi^{(i)}_{1,0}
=\psi^{(i),s=0}_1$、
$\psi^{(i)}_{1,1}
=\psi^{(i),s=1}_1$；
\item 满足条件 (4.6.26) 的
殆复结构双参数族
$\{J^{(i)}_{\rho,s}\}_{\rho,s}$；
此处须把其中的
$\{J^{s=j}_\rho\}_\rho$
等分别换成
$\{J^{(i),s=j}_\rho\}_\rho$
（$j=0,1$）；
\item 多值截面
$\mathfrak s_i^{\operatorname{twp}(s)}$、
$\mathfrak s_i^{
\operatorname{top}(\rho),\operatorname{twp}(s)}$；
\item 满足条件 (4.6.30) 的
$[0,1]\times L^{(i)\prime}$ 上
光滑奇异单形的可数集。
\end{itemize}
将这些同伦记作
\begin{equation}
\begin{aligned}
\mathfrak f^{
\psi^{(i)}_{1,s},J^{(i)}_{\rho,s},
\mathfrak s_i^{
\operatorname{top}(\rho),\operatorname{twp}(s)}}
\colon&
(C_1(L^{(i)};\Lambda_{0,\mathrm{nov}}),
\mathfrak m_k^{J_i,\mathfrak s_i})
\\
&\longrightarrow
(C([0,1]\times L^{(i)\prime})
\otimes\Lambda_{0,\mathrm{nov}},
\mathfrak m^{(i)}).
\end{aligned}
\tag{5.3.28}\label{eq:5-3-28}
\end{equation}
这里
$C([0,1]\times L^{(i)\prime})$
由上述可数集生成，
$\mathfrak m^{(i)}$
是利用 $\{J^{(i)}_{1,s}\}_s$
等构造的 $A_\infty$ 结构，
参见命题 4.6.37。

在此情形下，
我们将构造上述两套
滤过 $A_\infty$ 双模同态
$\phi_{s=0}$、$\phi_{s=1}$
之间的同伦。
为此取相容殆复结构的三参数族
$\{J_{\tau,t,s}\}_{\tau,t,s}$，
满足以下条件 (5.3.29)。
其中
$\{J^{(0)}_{\tau,s}\}_{\tau,s}$、
$\{J^{(1)}_{\tau,s}\}_{\tau,s}$
是用于构造 (5.3.28) 的双参数族；
$\{J_{\tau,t,s=0}\}_{\tau,t}$、
$\{J_{\tau,t,s=1}\}_{\tau,t}$
分别是构造 $\phi_{s=0}$、
$\phi_{s=1}$ 所用的双参数族。
\translatornote{原书实体页 323 在
(5.3.29) 的第三、四行及紧后的解释中，
将 $\psi$ 的首参数写作
$\chi^{-1}(\tau)$。
前文固定
$\chi\colon\mathbb R\to(0,1)$，
而此处 $\tau\in\mathbb R$，
故依先前移动边界公式统一为
$\chi(\tau)$；
见疑误 A-0301。}
\begin{equation}
\left\{
\begin{aligned}
J_{\tau,0,s}&=J^{(0)}_{\tau,s},\\
J_{\tau,1,s}&=J^{(1)}_{\tau,s},\\
J_{\tau,t,0}
&=\bigl(J_{\tau,t,s=0}\bigr)^{
\psi_{\chi(\tau),t,0}\circ
(\psi_{1,t,0})^{-1}},\\
J_{\tau,t,1}
&=\bigl(J_{\tau,t,s=1}\bigr)^{
\psi_{\chi(\tau),t,1}\circ
(\psi_{1,t,1})^{-1}},\\
J_{\tau,t,s}&=J_t
\qquad(\tau<-A),\\
J_{\tau,t,s}&=J_{\infty,t,s},
\quad\text{与 }\tau\text{ 无关}\quad(\tau>A).
\end{aligned}
\right.
\tag{5.3.29}\label{eq:5-3-29}
\end{equation}
其中 $A$ 充分大，
$\bigl(J_{\tau,t,s=i}\bigr)^{
\psi_{\chi(\tau),t,i}\circ
(\psi_{1,t,i})^{-1}}$
（$i=0,1$）
表示由微分同胚
$\psi_{\chi(\tau),t,i}
\circ(\psi_{1,t,i})^{-1}$
从 $J_{\tau,t,s=i}$
推送得到的殆复结构。
对每个固定的 $s$，
$\{J_{\infty,t,s}\}_t$
是一参数殆复结构族；
我们将用它在
$C([0,1]\times L^{(1)\prime})
\otimes\Lambda_{0,\mathrm{nov}}$
及
$C([0,1]\times L^{(0)\prime})
\otimes\Lambda_{0,\mathrm{nov}}$
分别作用的情形下，
构造滤过 $A_\infty$ 双模。

先定义承载这一双模结构的模。
回忆对每个
$h'\in\pi_0(L^{(0)\prime}\cap L^{(1)\prime})$
均须选择
$\mathcal X_3(R'_{h'})$；
还须证明结果不依赖这一选择。
设构造 $\phi_{s=0}$、$\phi_{s=1}$
分别使用
$\mathcal X_3^{s=0}(R'_{h'})$、
$\mathcal X_3^{s=1}(R'_{h'})$。
由下文引理 5.3.34，
可选定 $[0,1]\times R'_{h'}$
上光滑奇异链的可数集
$\mathcal X'([0,1]\times R'_{h'})$，
其系数取自 $\pi^*\Theta^-_{R'_{h'}}$，
并满足 (5.3.35)。
这里
$\pi\colon[0,1]\times R'_{h'}\to R'_{h'}$
是投影；不致混淆时，
仍以 $\Theta^-_{R'_{h'}}$
记 $\pi^*\Theta^-_{R'_{h'}}$。
该集合生成光滑奇异链复形
$S([0,1]\times R'_{h'};
\pi^*\Theta^-_{R'_{h'}})$
的子复形
$C([0,1]\times R'_{h'};
\pi^*\Theta^-_{R'_{h'}})$。
仿照第 3.7.5 小节，
定义
$\widehat{CF}_{BM,\mathrm{para}}
(L^{(1)\prime},L^{(0)\prime};
\Lambda_{\mathrm{nov}})$
为
\[
\bigl(
C([0,1]\times R'_{h'};
\pi^*\Theta^-_{R'_{h'}})
\otimes_{\mathbb Q}\Lambda_{h'}
\bigr)\otimes_{\mathbb Q}\Lambda_{\mathrm{nov}}
\]
的完备化，其中
\[
\Lambda_{h'}
=\operatorname{Span}_{\mathbb Q}
\bigl\{[h',w'_i]\mid
[h',w'_i]\in
\widetilde\pi_0(L^{(1)\prime}
\cap L^{(0)\prime}),\
h'\in\pi_0(L^{(1)\prime}
\cap L^{(0)\prime})\bigr\}.
\]
仿照第 3.7.4 小节的 (3.7.19)，
在这一完备化上定义等价关系 $\sim$。
置
\[
C_{\mathrm{para}}
(L^{(1)\prime},L^{(0)\prime};
\Lambda_{\mathrm{nov}})
=\widehat{CF}_{BM,\mathrm{para}}
(L^{(1)\prime},L^{(0)\prime};
\Lambda_{\mathrm{nov}})/\sim,
\]
\begin{equation}
\begin{aligned}
C_{\mathrm{para}}
(L^{(1)\prime},L^{(0)\prime};
\Lambda_{0,\mathrm{nov}})
=\;&C_{\mathrm{para}}
(L^{(1)\prime},L^{(0)\prime};
\Lambda_{\mathrm{nov}})\\
&\text{的非负能量部分。}
\end{aligned}
\tag{5.3.30}\label{eq:5-3-30}
\end{equation}

现在在 (5.3.30) 上定义滤过
$A_\infty$ 双模结构。
取
$[h'_1,w'_1],[h'_2,w'_2]
\in\widetilde\pi_0
(L^{(0)\prime}\cap L^{(1)\prime})$。
每个固定的 $s$ 下，
利用 $\{J_{\infty,t,s}\}_t$
定义模空间
\[
\mathcal M_{k_1,k_0}
\bigl(L^{(1)\prime},L^{(0)\prime};
[h'_1,w'_1],[h'_2,w'_2];
\{J_{\infty,t,s}\}_t\bigr),
\]
方法与 (3.7.29)、(3.7.56) 相同。
这里让 $t$ 变化，
考虑整个 $\{J_{\infty,t,s}\}_t$；
第 3.7.5 小节的 (3.7.54)
对此模空间采用了稍异的记号。
置
\begin{equation}
\begin{aligned}
&\mathcal M_{k_1,k_0}
\bigl(L^{(1)\prime},L^{(0)\prime};
[h'_1,w'_1],[h'_2,w'_2];
\{J_{\infty,t,s}\}_{t,s};
\operatorname{twp}(s)\bigr)\\
&\qquad=
\coprod_{s\in[0,1]}
\{s\}\times
\mathcal M_{k_1,k_0}
\bigl(L^{(1)\prime},L^{(0)\prime};
[h'_1,w'_1],[h'_2,w'_2];
\{J_{\infty,t,s}\}_t\bigr).
\end{aligned}
\tag{5.3.31}\label{eq:5-3-31}
\end{equation}
有评价映射
\begin{equation}
\begin{aligned}
ev={}&
(ev^{(1)}_1,\ldots,ev^{(1)}_{k_1};
ev_{-\infty};
ev^{(0)}_1,\ldots,ev^{(0)}_{k_0};
ev_{+\infty})\\
\colon&
\mathcal M_{k_1,k_0}
\bigl(L^{(1)\prime},L^{(0)\prime};
[h'_1,w'_1],[h'_2,w'_2];
\{J_{\infty,t,s}\}_{t,s};
\operatorname{twp}(s)\bigr)\\
\longrightarrow{}&
(L^{(1)\prime})^{k_1}
\times R'_{h'_1}
\times(L^{(0)\prime})^{k_0}
\times R'_{h'_2}.
\end{aligned}
\tag{5.3.32}\label{eq:5-3-32}
\end{equation}
与第 3.7.4、3.7.5 小节一样，
这一模空间具有 Kuranishi 结构，
其定向也可按第 3.7.5 小节的方法讨论。

取
$(|\mathcal P_i^{(0)}|,f_i^{(0)})
\in C([0,1]\times L^{(0)\prime})$、
$(|\mathcal P_i^{(1)}|,f_i^{(1)})
\in C([0,1]\times L^{(1)\prime})$
以及
$[S,f]\in
\mathcal X'([0,1]\times R'_{h'_1})$。
将 $f^{(j)}$ 和 $f$ 的坐标写作
\[
f^{(j)}(x)=
\bigl(f^{(j),s}(x),\bar f^{(j)}(x)\bigr)
\in[0,1]\times L^{(j)\prime},
\qquad
f(x)=\bigl(f^s(x),\bar f(x)\bigr).
\]
下一个定义中，简记
$[\mathcal P_i^{(j)}]$
代替
$[|\mathcal P_i^{(j)}|,f_i^{(j)}]$；
并记
$\mathcal P_i^{(j)}|_{s_0}
=(f_i^{(j),s})^{-1}(s_0)$，
$S|_{s_0}=(f^s)^{-1}(s_0)$。

\begin{definition}{5.3.33}\label{def:5-3-33}
对
$[S]\in\mathcal X'
([0,1]\times R'_{h'_1})$，
定义
\[
\mathfrak n_{k_1,k_0}
=\sum_{[h'_2,w'_2]}
\mathfrak n_{k_1,k_0;[h'_2,w'_2]},
\]
其中
\[
\begin{aligned}
&\mathfrak n_{k_1,k_0;[h'_2,w'_2]}
\bigl([\mathcal P_1^{(1)}],\ldots,
[\mathcal P_{k_1}^{(1)}],[S],
[\mathcal P_1^{(0)}],\ldots,
[\mathcal P_{k_0}^{(0)}]\bigr)\\
&\qquad=
T^{E(w'_2)-E(w'_1)}
e^{(\mu(w'_2)-\mu(w'_1))/2}[S'].
\end{aligned}
\]
这里 $[S']$ 如下定义。
考虑纤维积
\[
\begin{aligned}
|S'|={}&
\coprod_{s\in[0,1]}
\Bigl(\{s\}\times
\mathcal M_{k_1,k_0}
\bigl(L^{(1)\prime},L^{(0)\prime};
[h'_1,w'_1],[h'_2,w'_2];
\{J_{\infty,t,s}\}_t\bigr)\\
&\qquad\times_{
(L^{(1)\prime})^{k_1}
\times R'_{h'_1}
\times(L^{(0)\prime})^{k_0}}
\bigl(
\textstyle\prod_i\mathcal P_i^{(1)}|_s
\times S|_s
\times\prod_i\mathcal P_i^{(0)}|_s
\bigr)\Bigr).
\end{aligned}
\]
确切地说，这也可用
(5.3.31) 的参数化模空间
与 $\prod_i\mathcal P_i^{(1)}
\times S\times
\prod_i\mathcal P_i^{(0)}$
在
$([0,1]\times L^{(1)\prime})^{k_1}
\times([0,1]\times R'_{h'_1})
\times([0,1]\times L^{(0)\prime})^{k_0}$
上作纤维积来定义：
除评价映射 $ev$ 外，
所用映射还将
$(s,(\Sigma,\vec\tau,u))$ 送到
各个因子中的 $s\in[0,1]$，
而输入链各取自然投影。
空间 $S'$ 具有 Kuranishi 结构；
$[S']$ 是取多值截面并三角剖分
其零点集后得到的基本链，
系数取自 $\Theta^-_{R'_{h'_2}}$。
\end{definition}

\begin{lemma}{5.3.34}\label{lem:5-3-34}
可以选取
$[0,1]\times R'_{h'}$
上光滑奇异链的可数集
$\mathcal X'([0,1]\times R'_{h'})$
和定义 5.3.33 中的多值截面，
使其满足以下性质：
\begin{description}
\item[\textup{(5.3.35.1)}]\label{item:5-3-35-1}
定义 5.3.33 中的 $[S']$
可以分解为
$\mathcal X'([0,1]\times R'_{h'})$
中元素的和。
\item[\textup{(5.3.35.2)}]\label{item:5-3-35-2}
$\mathcal X'([0,1]\times R'_{h'})$
的每个元素都满足条件 4.6.30；
此处须将该条件中的 $L'$
换成 $R'_{h'}$。
\item[\textup{(5.3.35.3)}]\label{item:5-3-35-3}
对
$S\in\mathcal X'
([0,1]\times R'_{h'})$，
有
$S|_{s=0}\in\mathcal X_3^{s=0}(R'_{h'})$
及
$S|_{s=1}\in\mathcal X_3^{s=1}(R'_{h'})$。
这里
$*|_{s=s_0}
=\operatorname{Eval}_{s=s_0}(*)$
由 (4.6.32) 定义。
\item[\textup{(5.3.35.4)}]\label{item:5-3-35-4}
由这一可数集生成的链复形到
$S([0,1]\times R'_{h'};
\pi^*\Theta^-_{R'_{h'}})$
的包含映射是链同伦等价。
\end{description}
\end{lemma}

引理 5.3.34 的证明与其他推迟到第 7.2 节
证明的命题相似。
由 (5.3.35.1)，
$\mathfrak n_{k_1,k_0}$ 定义同态
\begin{equation}
\begin{aligned}
B_{k_1}C([0,1]\times L^{(1)\prime})
&\otimes
C_{\mathrm{para}}
(L^{(1)\prime},L^{(0)\prime};
\Lambda_{0,\mathrm{nov}})\\
&\otimes
B_{k_0}C([0,1]\times L^{(0)\prime})
\longrightarrow
C_{\mathrm{para}}
(L^{(1)\prime},L^{(0)\prime};
\Lambda_{0,\mathrm{nov}}).
\end{aligned}
\tag{5.3.36}\label{eq:5-3-36}
\end{equation}
直接检查可知，
这些同态构成一个滤过
$A_\infty$ 双模结构。
因此，链复形
$C_{\mathrm{para}}
(L^{(1)\prime},L^{(0)\prime};
\Lambda_{0,\mathrm{nov}})$
上已有所需的滤过双模结构。
又由 (5.3.35.2)、(5.3.35.3)，
存在滤过 $A_\infty$ 双模同态
\[
\operatorname{Eval}_{s=j}\colon
C_{\mathrm{para}}
(L^{(1)\prime},L^{(0)\prime};
\Lambda_{0,\mathrm{nov}})
\longrightarrow
C_3^{s=j}
(L^{(1)\prime},L^{(0)\prime};
\Lambda_{0,\mathrm{nov}})
\quad(j=0,1),
\]
它们分别立于
$(\operatorname{Eval}_{s=j},
\operatorname{Eval}_{s=j})$ 之上。
这里 $C_3^{s=j}$ 是由
$\mathcal X_3^{s=j}(R'_{h'})$
生成的子复形。
在引理 5.3.25 中，
我们都将它们记作 $C_3$；
此处标出 $j$，
以区分定义 $\phi_{s=0}$
和 $\phi_{s=1}$ 所用的两个复形。
由定理 5.2.35 和 (5.3.35.4)，
上述评价同态是同伦等价；
这里系数为
$\Lambda_{0,\mathrm{nov}}$，
故可以应用该定理。

下面构造弱滤过
$A_\infty$ 双模同态
\[
\begin{aligned}
\phi^{
\psi_{\rho,s}^{(*)},J_{\tau,t,s},
\mathfrak s^{\operatorname{top}(\tau),
\operatorname{twp}(s)}}\colon
&C(L^{(1)},L^{(0)};
\Lambda_{\mathrm{nov}})\\
&\longrightarrow
C_{\mathrm{para}}
(L^{(1)\prime},L^{(0)\prime};
\Lambda_{\mathrm{nov}}),
\end{aligned}
\]
它立于一对滤过代数同态
\[
\left(
\mathfrak f^{
\psi^{(1)}_{1,s},J^{(1)}_{\rho,s},
\mathfrak s^{\operatorname{top}(\rho),
\operatorname{twp}(s)}_1},
\mathfrak f^{
\psi^{(0)}_{1,s},J^{(0)}_{\rho,s},
\mathfrak s^{\operatorname{top}(\rho),
\operatorname{twp}(s)}_0}
\right)
\]
之上，并满足
\begin{equation}
\operatorname{Eval}_{s=j}
\circ
\phi^{
\psi_{\rho,s}^{(*)},J_{\tau,t,s},
\mathfrak s^{\operatorname{top}(\tau),
\operatorname{twp}(s)}}
=\phi_{s=j}^{
\psi^{(0),s=j},\psi^{(1),s=j},
\{J_{\tau,t,s=j}\}_{\tau,t},
\mathfrak s_{s=j}}
\quad(j=0,1).
\tag{5.3.37}\label{eq:5-3-37}
\end{equation}
这里回忆
$\rho\in[0,1]$ 与
$\tau\in\mathbb R$
由 $\rho=\chi(\tau)$ 联系。
为排出下列交换图，
暂记
\[
\begin{aligned}
D&=C(L^{(1)},L^{(0)};
\Lambda_{\mathrm{nov}}),\\
D_{\mathrm{para}}
&=C_{\mathrm{para}}
(L^{(1)\prime},L^{(0)\prime};
\Lambda_{\mathrm{nov}}),\\
D_3^j&=C_3^{s=j}
(L^{(1)\prime},L^{(0)\prime};
\Lambda_{\mathrm{nov}})
\quad(j=0,1).
\end{aligned}
\]
\begin{center}
\[
\begin{array}{ccc}
D&\xrightarrow{\ \phi_{s=0}\ }&D_3^0\\[2mm]
\big\Vert&&\Big\uparrow\vcenter{\hbox{\scriptsize
$\operatorname{Eval}_{s=0}$}}\\[2mm]
D&\longrightarrow&D_{\mathrm{para}}\\[2mm]
\big\Vert&&\Big\downarrow\vcenter{\hbox{\scriptsize
$\operatorname{Eval}_{s=1}$}}\\[2mm]
D&\xrightarrow{\ \phi_{s=1}\ }&D_3^1 .
\end{array}
\]
\renewcommand{\thefigure}{5.3.3}
\refstepcounter{figure}
{\normalfont\bfseries 图式~\thefigure}\label{diag:5-3-3}
\end{center}

这一构造是上一小节构造的参数族版本。
对每个 $s$，
考虑相容殆复结构的双参数族
$\{J_{\tau,t,s}\}_{\tau,t}$
及哈密顿同痕族
$\{\psi^{(i)}_{\rho,s}\}_{\rho,s}$
（$i=0,1$）。
如上一小节，
它们给出模空间
\[
\mathcal M_{k_1,k_0}
\bigl(
L^{(1)},L^{(0)};
\{J_{\tau,s}\}_\tau;
[h,w],[h',w'];
\operatorname{top}(\tau)
\bigr)
\]
（每个 $s$ 各一个）。
这里以
$\{J_{\tau,s}\}_\tau$
表示固定 $s$ 后的
$t$-方向殆复结构路径族
$\{J_{\tau,t,s}\}_{\tau,t}$。
再置
\[
\begin{aligned}
&\mathcal M_{k_1,k_0}
\bigl(
L^{(1)},L^{(0)};
\{J_{\tau,s}\}_{\tau,s};
[h,w],[h',w'];
\operatorname{top}(\tau),
\operatorname{twp}(s)
\bigr)\\
&\qquad=
\coprod_{s\in[0,1]}\{s\}\times
\mathcal M_{k_1,k_0}
\bigl(
L^{(1)},L^{(0)};
\{J_{\tau,s}\}_{\tau};
[h,w],[h',w'];
\operatorname{top}(\tau)
\bigr).
\end{aligned}
\]
利用此空间定义
\[
\phi^{
\psi_{\rho,s}^{(*)},J_{\tau,t,s},
\mathfrak s^{\operatorname{top}(\tau),
\operatorname{twp}(s)}}
=
\bigl\{
\phi_{k_1,k_0}^{
\psi_{\rho,s}^{(*)},J_{\tau,t,s},
\mathfrak s^{\operatorname{top}(\tau),
\operatorname{twp}(s)}}
\bigr\}_{k_1,k_0}.
\]
下式中的各 $\phi$ 分量
均带有这个相同的上标。
沿用 $\mathcal P_i^{(j)}$、$S$
表示输入链的记法；
此处它们取自源端的
$L^{(j)}$、$R_h$，而非
定义 5.3.33 的目标端参数化复形。
令
\[
\phi_{k_1,k_0}
=\sum_{[h',w']}
\phi_{k_1,k_0;[h',w']},
\]
其中
\[
\begin{aligned}
&\phi_{k_1,k_0;[h',w']}
\bigl(
[\mathcal P_1^{(1)}],\ldots,
[\mathcal P_{k_1}^{(1)}],[S],
[\mathcal P_1^{(0)}],\ldots,
[\mathcal P_{k_0}^{(0)}]
\bigr)\\
&\qquad=
T^{E(w')-E(w)}
e^{(\mu(w')-\mu(w))/2}[S'].
\end{aligned}
\]
这里 $S'$ 是纤维积
\[
\begin{aligned}
S'=\coprod_{s\in[0,1]}
\Bigl(\{s\}\times
&\mathcal M_{k_1,k_0}
\bigl(
L^{(1)},L^{(0)};
\{J_{\tau,s}\}_{\tau,s};
[h,w],[h',w'];
\operatorname{top}(\tau),
\operatorname{twp}(s)
\bigr)\\
&\times_{
(L^{(1)})^{k_1}\times R_h
\times(L^{(0)})^{k_0}}
\bigl(
\textstyle\prod_i\mathcal P_i^{(1)}
\times S
\times\prod_i\mathcal P_i^{(0)}
\bigr)\Bigr).
\end{aligned}
\]
这个空间具有 Kuranishi 结构；
$[S']$ 是选取多值截面
并三角剖分其零点集得到的基本链。
可一并选取多值截面及
$\mathcal X'([0,1]\times R'_{h'})$，
使 $[S']$ 分解为该集合中的元素，
从而上述 $\phi$ 得到定义。

与上一小节相同的论证表明，
所得映射是
$A_\infty$ 双模同态。
注意 $s=0,1$ 时的模空间
分别与定义
$\phi_{s=0}$、
$\phi_{s=1}$
所用的模空间一致；
因而 (5.3.37) 成立。

\begin{lemma}{5.3.38}\label{lem:5-3-38}
存在滤过 $A_\infty$ 双模同态
\[
\operatorname{Incl}\colon
C_2(L^{(1)\prime},L^{(0)\prime};
\Lambda_{0,\mathrm{nov}})_{\mathrm{para}}
\longrightarrow
C_{\mathrm{para}}
(L^{(1)\prime},L^{(0)\prime};
\Lambda_{0,\mathrm{nov}}),
\]
它立于
$(\operatorname{Incl}^{L^{(1)}},
\operatorname{Incl}^{L^{(0)}})$ 之上，
其中
\[
\begin{aligned}
\operatorname{Incl}^{L^{(i)}}\colon&
\bigl(C_2(L^{(i)\prime};
\Lambda_{0,\mathrm{nov}}),
\mathfrak m_k^{J'_i,\mathfrak s'_i}\bigr)\\
&\longrightarrow
\bigl(
C_2([0,1]\times L^{(i)\prime};
\Lambda_{0,\mathrm{nov}}),\mathfrak m
\bigr)
\end{aligned}
\]
如 (4.6.38.8) 所定义。
复合映射
\[
\operatorname{Eval}_{s=s_0}
\circ\operatorname{Incl}\colon
C_2(L^{(1)\prime},L^{(0)\prime};
\Lambda_{0,\mathrm{nov}})
\longrightarrow
C_3^{s=s_0}
(L^{(1)\prime},L^{(0)\prime};
\Lambda_{0,\mathrm{nov}})
\]
等于包含映射。
\end{lemma}

\begin{proof}
仿照 (4.6.38.8)，
利用棱柱型单纯剖分定义
\[
\operatorname{Incl}_{1,\beta_0}\colon
C_2(L^{(i)\prime};\mathbb Q)
\longrightarrow
C_2([0,1]\times L^{(i)\prime};
\mathbb Q).
\]
再仿照 (4.6.38.8)、
引理 5.3.34
及命题 7.2.162 的情形，
定义 $\operatorname{Incl}_{k,\beta}$。
细节与滤过 $A_\infty$ 代数的情形相似，
从略。
\end{proof}

引理 5.3.38 中的
$\operatorname{Incl}$ 定义在
$\Lambda_{0,\mathrm{nov}}$ 上，
故可由定理 5.2.35 知它是同伦等价。
于是可将图式 5.3.3 扩充如下。
为排图，记
$D_2=C_2(L^{(1)\prime},L^{(0)\prime};
\Lambda_{\mathrm{nov}})$；
图中各行的 $D_2$
表示引理 5.3.38 中相应的双模，
省略参数族下标及系数记号。
\begin{center}
\[
\begin{array}{ccccc}
D&\xrightarrow{\ \phi_{s=0}\ }&
D_3^0&\xleftarrow{\ \mathrm{incl}\ }&D_2\\[2mm]
\big\Vert&&
\Big\uparrow\vcenter{\hbox{\scriptsize
$\operatorname{Eval}_{s=0}$}}
&&\big\Vert\\[2mm]
D&\longrightarrow&
D_{\mathrm{para}}
&\xleftarrow{\ \operatorname{Incl}\ }&D_2\\[2mm]
\big\Vert&&
\Big\downarrow\vcenter{\hbox{\scriptsize
$\operatorname{Eval}_{s=1}$}}
&&\big\Vert\\[2mm]
D&\xrightarrow{\ \phi_{s=1}\ }&
D_3^1&\xleftarrow{\ \mathrm{incl}\ }&D_2 .
\end{array}
\]
\renewcommand{\thefigure}{5.3.4}
\refstepcounter{figure}
{\normalfont\bfseries 图式~\thefigure}\label{diag:5-3-4}
\end{center}
图式右侧的三个横向箭头
都是同伦等价。
现在仿照引理 4.6.43 的证明，
即可证明定理 5.3.27。
\end{proof}
